\section{Proposed Methodology}

This research adopts a cross-scale and integrative methodology, establishing a bidirectional coupling framework that links neighbourhood microclimates with building energy performance. Sydney is selected as a case study for empirical validation because it provides a suitable warm-temperate coastal context with increasing summer heat stress, strong intra-urban microclimate variation, high-density inner-city neighbourhoods, established net-zero building policies, and relatively rich public datasets on climate, urban form, building stock, and energy consumption. While the proposed framework is intended to be broadly applicable, Sydney provides a suitable context to test its effectiveness and refine methodological strategies. The methodology follows four main steps: identifying and quantifying key factors, analysing interaction mechanisms, developing the coupling framework, and validating the framework, allowing systematic investigation of neighbourhood--building energy interactions.

\subsection{Data Collection and Key Factor Identification}

To identify key factors linking neighbourhood microclimates and building energy performance, this study will collect multi-source environmental, building, and human behaviour data. Official meteorological records and reanalysis data, such as ERA5-Land and national or provincial meteorological station data, will provide baseline climate information. On-site microclimate measurements will be conducted in 10 representative neighbourhoods differing in density, greenery, and morphology, capturing key parameters such as air temperature, relative humidity, and thermal indices. Questionnaires will collect at least 200 residents' thermal comfort perceptions, behavioural responses, and building energy use patterns. Supplementary aggregated social media data may also be used to capture extreme thermal events and community energy behaviour trends.

Collected datasets will be integrated through statistical correlation, regression modelling, and spatial analysis. Specifically, meteorological and morphological indicators such as SVF, NDVI, and H/W ratio will be aligned with energy consumption and behavioural survey data to test associations. Multivariate regression and variance decomposition will quantify the relative contributions of environmental, morphological, and behavioural factors. This integration will provide an evidence-based foundation for identifying the most influential determinants of the performance relationships among microclimate, thermal comfort perception, behaviour, and energy consumption.

\subsection{Interaction Mechanism Analysis}

To reveal the causal linkages between neighbourhood microclimates, thermal perception, human behaviour, and building energy consumption, this stage focuses on quantitative mechanism analysis based on the integrated datasets. The multi-source data on meteorology, urban morphology, and energy use will first be preprocessed through outlier detection and missing-value imputation to ensure accuracy and consistency. Correlation and regression analyses will then be applied to identify the statistical associations between microclimate parameters and energy consumption while controlling for building attributes and temporal effects.

Building on thermal comfort and behavioural theories, logistic regression will be used to estimate the probability of air-conditioning usage under varying UTCI conditions, while ordinal logistic regression will model the relationship between thermal indices and occupants' thermal sensation or acceptability levels. A path model will then be developed to link UTCI, thermal perception, adaptive behaviour, and energy consumption, allowing the identification of both direct and indirect effects and the mediating role of human behaviour in the microclimate--energy relationship. Finally, spatial analysis will be conducted to examine the spatial heterogeneity of these mechanisms across different neighbourhoods, revealing how morphological characteristics modulate behavioural and energy responses to microclimate variations.

\subsection{Development of the Bidirectional Coupling Framework}

This stage develops a bidirectional coupling framework that integrates microclimate simulation, building energy modelling, and resident behavioural dynamics into a unified system. ENVI-met and UWG will be employed to simulate microclimate conditions at neighbourhood and city scales, providing localised boundary conditions for subsequent energy simulations. These simulated microclimate data will be imported into EnergyPlus or CityBES to calculate building energy performance under different environmental scenarios.

To capture human adaptive effects, a resident behavioural model will be embedded into the energy simulation through a Python interface, enabling dynamic feedback between thermal perception, behavioural response, and energy consumption. Furthermore, a simplified bidirectional coupling will be established by feeding building waste heat outputs back into the microclimate simulation, thereby forming an interactive human--environment--building system. Scenario simulations will be conducted to test the impacts of urban design variables on outdoor thermal comfort and indoor energy demand. Model calibration and optimisation using measured data will ensure the reliability, stability, and applicability of the coupled framework.

\subsection{Framework Validation and Application Potential Assessment}

The final stage validates the proposed framework and evaluates its potential for transferability to other cities. After calibration using empirical measurements in Sydney, the framework will be applied to new urban contexts to test its adaptability and robustness. Only minimal secondary datasets are required for model transfer, including official meteorological records, remote sensing or GIS-derived morphological indicators such as SVF, NDVI, and albedo, and aggregated household energy statistics. Behavioural information can be approximated from social media data or existing survey databases, reducing the need for large-scale fieldwork.

During localisation, input parameters such as climatic distributions, morphological features, and behavioural coefficients will be updated to reflect the characteristics of the new city, enabling rapid adaptation without extensive retraining. Model validation will involve comparing simulated results with limited measured meteorological and energy data to ensure that the outputs are of reasonable magnitude and follow expected temporal patterns. Finally, scenario analyses will be conducted to derive localised energy--comfort optimisation thresholds, demonstrating the framework's applicability for guiding climate-responsive and energy-efficient urban design strategies.
